\chapter{Functions of a Real Variable}\label{ch:b2:realfun}

Before analysis moves to functions of functions (\cref{ch:b2:funcseq})
it pays to know the one-variable landscape in finer detail than
Year 1 required: how discontinuous a monotone function can be, how
regular a convex function must be, and what special properties
derivatives enjoy (Darboux). These structural results are short,
sharp, and beloved of examiners.

\section{Monotone functions}

\begin{theorem}[Regularity of monotone functions]\label{thm:b2:realfun:monotone}
Let $f \colon I \to \R$ be increasing on an interval.
\begin{enumerate}
  \item At every interior point $a$, the one-sided limits exist:
        \[
        f(a^-) = \sup_{x < a} f(x) \;\leq\; f(a) \;\leq\;
        f(a^+) = \inf_{x > a} f(x) ;
        \]
        every discontinuity is a \emph{jump}.
  \item The set of discontinuities of $f$ is at most
        countable.\index{monotone function!discontinuities}
\end{enumerate}
\end{theorem}

\begin{proof}
(1) The set $\{f(x) : x < a\}$ is nonempty, bounded above by $f(a)$:
its supremum $s$ satisfies $f(x) \to s$ as $x \to a^-$ (given
$\varepsilon$, some $f(x_0) > s - \varepsilon$, and monotonicity
traps $f(x) \in \intoc{s - \varepsilon}{s}$ for $x \in
\intoo{x_0}{a}$). Symmetrically on the right.

(2) To each discontinuity $a$ attach the nonempty open interval
$J_a = \intoo{f(a^-)}{f(a^+)}$ (a genuine jump). For $a < b$
discontinuities, $J_a$ and $J_b$ are disjoint: $f(a^+) \leq f(c)
\leq f(b^-)$ for any $c$ between. Each $J_a$ contains a rational;
distinct discontinuities get distinct rationals: an injection of the
discontinuity set into $\Q$, which is countable
(\cref{prop:b2:structures:countablestable}).
\end{proof}

\begin{example}
The bound is sharp: fix an enumeration $(r_n)$ of $\Q \cap
\intoo{0}{1}$ and set $f(x) = \sum_{n : r_n \leq x} 2^{-n}$
(a summable-family definition,
\cref{def:b2:series:summable}). Then $f$ is increasing on
$\intcc{0}{1}$ and discontinuous exactly at every rational of
$\intoo{0}{1}$ (jump $2^{-n}$ at $r_n$): a monotone function
\emph{can} be discontinuous on a dense countable set.
\end{example}

\section{Convex functions}

\begin{lemma}[Slope inequality]\label{lem:b2:realfun:slopes}
Let $f$ be convex on $I$ and $x < y < z$ in $I$. Then
\[
\frac{f(y) - f(x)}{y - x}
\;\leq\; \frac{f(z) - f(x)}{z - x}
\;\leq\; \frac{f(z) - f(y)}{z - y} :
\]
slopes of chords increase in both endpoints.\index{convex
function!slope inequality}
\end{lemma}

\begin{proof}
Write $y = \frac{z - y}{z - x}\,x + \frac{y - x}{z - x}\,z$ (a
convex combination) and expand $f(y) \leq \frac{z-y}{z-x}f(x) +
\frac{y-x}{z-x}f(z)$: rearranging once gives the left inequality,
rearranging the other way the right one.
\end{proof}

\begin{theorem}[Regularity of convex functions]\label{thm:b2:realfun:convexreg}
Let $f$ be convex on an interval $I$.
\begin{enumerate}
  \item At every interior point, $f$ has finite one-sided
        derivatives $f'_g \leq f'_d$; both are increasing functions
        of the point; in particular $f$ is continuous on the
        interior of $I$ (but possibly not at endpoints).
  \item $f$ lies above each \emph{support line}: for $a$ interior
        and any $m \in \intcc{f'_g(a)}{f'_d(a)}$,
        \[
        f(x) \geq f(a) + m(x - a) \qquad (x \in I).
        \]
  \item (Jensen, weighted) For $x_i \in I$ and weights $\lambda_i
        \geq 0$, $\sum\lambda_i = 1$:
        \[
        f\Bigl(\sum_i \lambda_i x_i\Bigr) \leq \sum_i \lambda_i
        f(x_i) .\index{Jensen's inequality}
        \]
\end{enumerate}
\end{theorem}

\begin{proof}
(1) Fix $a$ interior. By \cref{lem:b2:realfun:slopes}, the slope
$\tau(h) = \frac{f(a + h) - f(a)}{h}$ is an increasing function of
$h$ (on both sides, and $\tau(h_-) \leq \tau(h_+)$ for $h_- < 0 <
h_+$). Hence $\tau$ has a finite limit as $h \to 0^-$ (increasing,
bounded above by any right slope) --- this is $f'_g(a)$ --- and as
$h \to 0^+$ ($f'_d(a)$), with $f'_g(a) \leq f'_d(a)$. Finite
one-sided derivatives force continuity at $a$. Monotonicity in the
point: for $a < b$ interior, $f'_d(a) \leq \frac{f(b) - f(a)}{b -
a} \leq f'_g(b)$, again by the slope inequality.

(2) For $x > a$: $\frac{f(x) - f(a)}{x - a} \geq f'_d(a) \geq m$;
for $x < a$: $\frac{f(a) - f(x)}{a - x} \leq f'_g(a) \leq m$. Both
rearrange to the claim.

(3) Induction on the number of points exactly as in the Year 1
volume (the two-point case is the definition) --- or in one stroke:
apply (2) at $a = \sum\lambda_i x_i$ and average the support-line
inequalities at the points $x_i$ with weights $\lambda_i$: $\sum_i
\lambda_i f(x_i) \geq f(a) + m\sum_i\lambda_i(x_i - a) = f(a)$.
\end{proof}

\begin{example}[Endpoint discontinuity]
On $\intcc{0}{1}$, the function $f(0) = 1$, $f(x) = 0$ for $x > 0$
is convex but discontinuous at the endpoint $0$: statement (1) is
sharp.
\end{example}

\begin{example}[Power mean inequality]\label{ex:b2:realfun:powermeans}
For $0 < p < q$ and positive $x_i$ with weights $\lambda_i$
summing to $1$, applying Jensen to the convex $t \mapsto t^{q/p}$
at the points $x_i^p$:
\[
\Bigl(\sum \lambda_i x_i^{p}\Bigr)^{1/p}
\leq \Bigl(\sum \lambda_i x_i^{q}\Bigr)^{1/q} :
\]
power means increase with the exponent --- containing AM--QM, and,
in the limit $p \to 0$ (\cref{exo:b2:realfun:6}), the AM--GM
inequality once more.
\end{example}

\section{The Darboux property}

\begin{theorem}[Darboux]\label{thm:b2:realfun:darboux}
Let $f$ be differentiable on an interval $I$. Then $f'$ takes every
value between any two of its values --- even though $f'$ need not be
continuous.\index{Darboux's theorem}
\end{theorem}

\begin{proof}
Let $a < b$ in $I$ and $v$ strictly between $f'(a)$ and $f'(b)$,
say $f'(a) < v < f'(b)$. The function $g(x) = f(x) - vx$ is
differentiable with $g'(a) < 0 < g'(b)$: its minimum on
$\intcc{a}{b}$ (attained: continuity on a compact) is not at $a$
(just after $a$, $g$ decreases below $g(a)$) nor at $b$ (just
before $b$, $g$ is below $g(b)$): it is interior, and there $g'(c)
= 0$, i.e.\ $f'(c) = v$. (This was a Year 1 starred exercise; its
place in the theory is here.)
\end{proof}

\begin{corollary}\label{cor:b2:realfun:nojumps}
A derivative has no jump discontinuities: if $f'(a^-)$ and
$f'(a^+)$ exist, they equal $f'(a)$. The discontinuities of a
derivative are always of oscillation type ($x^2\sin\frac1x$'s
derivative at $0$, Year 1 volume).
\end{corollary}

\begin{proof}
If $f'(a^+) = \lim_{x\to a^+} f'(x)$ exists and differs from
$f'(a)$, values strictly between them would be skipped by $f'$ on
a right neighborhood --- contradicting Darboux on intervals
$\intcc{a}{a + h}$. (Alternatively: the mean value theorem forces
$f'(a) = \lim_{h\to0^+} \frac{f(a+h)-f(a)}{h} = f'(a^+)$, the
difference quotient being an $f'$-value at an intermediate point.)
Same on the left.
\end{proof}

\section{Exercises}

\begin{exercise}[$\star$]\label{exo:b2:realfun:1}
Determine the discontinuity sets and the jump sizes: $\lfloor x
\rfloor$; $\;x - \lfloor x\rfloor$; $\;\lfloor x \rfloor + \sqrt{x
- \lfloor x\rfloor}$; the function of the example following
\cref{thm:b2:realfun:monotone} restricted to dyadic rationals
$r_n$.
\end{exercise}

\begin{exercise}[$\star$]\label{exo:b2:realfun:2}
Prove that an increasing function $f \colon I \to \R$ with the
intermediate value property (its image of any subinterval is an
interval) is continuous.
\end{exercise}

\begin{exercise}[$\star$]\label{exo:b2:realfun:3}
Which of the following are convex on their domain? $x \mapsto
x\ln x$ ($x > 0$); $\;x \mapsto \ln(1 + \eu^x)$; $\;x \mapsto
\sqrt{1 + x^2}$; $\;x \mapsto x^3$.
\end{exercise}

\begin{exercise}[$\star\star$]\label{exo:b2:realfun:4}
Let $f$ be convex on $\R$ and bounded above. Prove that $f$ is
constant. \emph{(If $f(a) \neq f(b)$, the slope inequality
propagates the nonzero chord slope: beyond the point with the larger
value, $f$ grows at least linearly --- contradicting boundedness.
Treat both signs of the slope.)} Deduce that a convex function on
$\R$ with an asymptote at both ends is affine.
\end{exercise}

\begin{exercise}[$\star\star$]\label{exo:b2:realfun:5}
Let $f$ be differentiable on $I$ with $f'$ monotone. Prove that
$f'$ is continuous \emph{(combine
\cref{thm:b2:realfun:monotone} and
\cref{cor:b2:realfun:nojumps})}.
\end{exercise}

\begin{exercise}[$\star\star$]\label{exo:b2:realfun:6}
(Geometric mean as a limit) For positive $x_i$ and weights
$\lambda_i$ summing to $1$, prove
\[
\lim_{p \to 0^+} \Bigl(\sum_i \lambda_i x_i^p\Bigr)^{1/p}
= \prod_i x_i^{\lambda_i} ,
\]
via $x_i^p = \eu^{p\ln x_i} = 1 + p\ln x_i + O(p^2)$, and deduce
the weighted AM--GM inequality from
\cref{ex:b2:realfun:powermeans}.
\end{exercise}

\begin{exercise}[$\star\star$]\label{exo:b2:realfun:7}
(Entropy inequality) Using strict convexity of $t \mapsto t\ln t$,
prove that for positive $p_i, q_i$ with $\sum p_i = \sum q_i = 1$:
\[
\sum_i p_i \ln\frac{p_i}{q_i} \geq 0 ,
\]
with equality iff $p = q$. \emph{(Write the left side as $\sum q_i\,
\varphi\bigl(\frac{p_i}{q_i}\bigr)$ with $\varphi(t) = t\ln t$ and
apply Jensen with weights $q_i$.)}
\end{exercise}

\begin{exercise}[$\star\star\star$]\label{exo:b2:realfun:8}
(Midpoint convexity) $f \colon I \to \R$ is \emph{midpoint convex}
when $f\bigl(\frac{x+y}{2}\bigr) \leq \frac{f(x) + f(y)}{2}$
always. Prove that a \emph{continuous} midpoint convex function is
convex. \emph{(Establish the convexity inequality for dyadic
weights $\frac{k}{2^m}$ by induction on $m$, then pass to the limit
using density and continuity.)}
\end{exercise}

\begin{exercise}[$\star\star\star$]\label{exo:b2:realfun:9}
Let $f$ be convex on $\intco{0}{+\infty}$ with $f(0) \leq 0$. Prove
that $x \mapsto \frac{f(x)}{x}$ is increasing on
$\intoo{0}{+\infty}$, and deduce that for convex $f$ with $f(0) =
0$: $f(x + y) \geq f(x) + f(y)$ for $x, y \geq 0$
(superadditivity).
\end{exercise}
